\clearpage
\fancyhead[L]{\small BOUVARD. JUPITER UND SATURN.}
\SourceLeaf{513}
\begin{center}\small Göttingische gelehrte Anzeigen. Stück 177. Seite 1761..1764. 1808 November 5.\\[3pt]\rule{.72\linewidth}{.3pt}\end{center}

Chez Courcier: \emph{Tables astronomiques publiées par le bureau des longitudes de France. Nouvelles tables de Jupiter et de Saturne calculées d'après la théorie de M. }\textsc{Laplace}, \emph{et suivant la division decimale de l'angle droit; par M.} \textsc{Bouvard}. 1808. 4. Paris.

Diese Tafeln, welche als die Fortsetzung der im J. 1806 herausgegebenen und bereits im 95. Stück dieser Blätter von 1807 angezeigten Sonnen- und Mondtafeln angesehen werden müssen, sind die köstliche Frucht von den mehr als 20 Jahre hindurch fortgesetzten und immer mehr vervollkommneten Untersuchungen \textsc{Laplace}'s über die Störungen, wodurch die Bewegungen der beiden grössten Planeten unsers Sonnensystems so verwickelt werden. Mit welchem glücklichen Erfolge dieser grosse Geometer hier, so wie in allen andern Theilen der physischen Astronomie, die er zum Gegenstande seiner Forschungen machte, alle Schwierigkeiten besiegt, und die widerspenstigen Bewegungen dem Calcul unterworfen hat, ist bekannt; ihm verdanken wir es, dass wir jetzt diese Bewegung fast völlig eben so genau berechnen als beobachten können. Die auf die ersten \textsc{Laplace}'schen Untersuchungen gegründeten Jupiters- und Saturnstafeln, welche \textsc{Delambre} berechnet hatte, wichen schon selten um eine halbe Minute vom Himmel ab; allein die seitdem von \textsc{Laplace} noch viel weiter getriebene Berechnung der Störungen, und die neue Bestimmung der elliptischen Elemente, welche von \textsc{Bouvard}, mit Zuziehung von jenen, aus der sorgfältigsten Discussion aller seit einem halben Jahrhundert beobachteten Oppositionen entwickelt sind, haben eine drei Mal so grosse Uebereinstimmung zur Folge gehabt. Die letzten Resultate dieser Arbeit, nemlich die in Formeln gebrachten Bewegungen des Jupiter und Saturn, sind schon in dem 8. Theile der Mécanique céleste S. 337 bekannt gemacht: im gegenwärtigen Werke finden wir nun dieselben in Tafeln gebracht, über deren Einrichtung wir nur noch wenig hinzu zu setzen haben, da sie in den meisten Stücken mit derjenigen übereinstimmt, die schon aus den \textsc{Delambre}'schen Sonnentafeln bekannt ist. Die Zeiten werden nicht vom Mittage, sondern von der Mitternacht an gezählt; die Anomalien von der Sonnennähe, und die Gleichungen sind dadurch alle positiv gemacht, dass man zu jeder eine beständige Grösse hinzusetzte, und die Summe aller dieser Vermehrungen wieder von der Mittelpunktsgleichung (die, wo es nöthig war, durch Hinzufügung von 400 Decimalgraden positiv gemacht wurde), dem elliptischen Radius Vector, und dem statt der Breite eingeführten Abstande vom Nordpol wieder abzog. Alle diese Einrichtungen (die neue, ohne irgend einen sichtbaren Nutzen angenommene, Art, die Zeit zu zählen, abgerechnet) haben längst den Beifall der Astronomen erhalten. Die Befolgung der Decimaleintheilung des Quadranten in diesen 'Tafeln soll ein Versuch sein die Astronomen nach und nach daran zu gewöhnen: so lange indess zugleich Sonnen- und Sinustafeln nach der alten Eintheilung gebraucht werden müssen, ist dies mehr eine Unbequemlichkeit, als Erleichterung.

Die Einleitung, welche den Tafeln vorgesetzt ist, gibt zuvörderst die Beschreibung der einzelnen Tafeln; hiernächst die Formeln, wonach sie berechnet sind, und die von den in der Mécanique céleste a. a. O. gegebenen weiter nicht verschieden sind, als dass man alle Epochen und veränderliche Coëfficienten von 1750 auf 1800 reducirt hat, endlich die vollständige Berechnung einer am 2. April 1806 gemachten Jupitersbeobachtung, wo die Tafeln um $20^{\prime\prime}0$ in der Länge, und um $19^{\prime\prime}7$ in der Breite abweichen ($6^{\prime\prime}8$ und $6^{\prime\prime}4$ nach der Sexagesimaleintheilung). Dies ist der Fehler für den geocentrischen Ort: es wird hierauf noch gezeigt, wie man daraus den Fehler des heliocentrischen ableiten kann; da inzwischen hier der berechnete Radius Vector als vollkommen genau angesehen werden muss, so hätte die Erinnerung hier nicht fehlen dürfen, dass man dieses Verfahren nur dann mit einiger Sicherheit anwen\SourceLeaf{514}den darf, wo ein mässiger Fehler im Radius Vector nur einen sehr geringen Einfluss auf die Länge und Breite hat, also nur bei den entfernten Planeten oder in der Nähe der Opposition, so wie bei mässigen Neigungen, oder in der Nähe der Knoten.

Da diese Jupiters- und Saturnstafeln lediglich auf die neuern Beobachtungen gegründet sind, so ist es interessant, zu sehen, mit welcher Genauigkeit ältere Beobachtungen dadurch dargestellt werden. Vorzüglich merkwürdig ist in dieser Rücksicht die am 31. October 1007 zu Cairo von \textsc{Ibn Junis} beobachtete Zusammenkunft jener beiden Planeten; nach \textsc{Bouvard}'s Rechnung findet sich hier zwischen den Tafeln und der Beobachtung ein Unterschied von $11'\,25''$ in der Länge, und von $6'\,56''$ in der Breite (Sexagesimal-Eintheilung), welcher allerdings den möglichen Beobachtungsfehler kaum übersteigt. Inzwischen bemerken wir hier noch, dass nach Erscheinung der gegenwärtigen Tafeln sich bei der grossen Gleichung, sowohl für den Jupiter als für den Saturn, ein kleiner Fehler gefunden hat, indem ein Glied mit unrechtem Zeichen genommen war. Diese Veränderung in der grossen Gleichung macht aber zugleich eine kleine Aenderung bei den Epochen und mittlern Bewegungen nothwendig: das Resultat davon ist, dass zu $q^{\mathrm{IV}}$ (Mécan. cél. Band 4 S. 338) noch hinzugefügt werden muss:
\[
\begin{aligned}
&51''98 + t.0''4156\\
&{}- (73''58 - t.0''0103)\sin(5n^{\mathrm{V}}t - 2n^{\mathrm{IV}}t + 5\varepsilon^{\mathrm{V}} - 2\varepsilon^{\mathrm{IV}})\\
&{}+ (47''63 + t.0''0287)\cos(5n^{\mathrm{V}}t - 2n^{\mathrm{IV}}t + 5\varepsilon^{\mathrm{V}} - 2\varepsilon^{\mathrm{IV}})
\end{aligned}
\]
und zu $q^{\mathrm{V}}$ (ebendas.)
\[
\begin{aligned}
&-127''13 - t.1''0212\\
&{}+ (179''952 - t.0''025192)\sin(5n^{\mathrm{V}}t - 2n^{\mathrm{IV}}t + 5\varepsilon^{\mathrm{V}} - 2\varepsilon^{\mathrm{IV}})\\
&{}- (116''541 + t.0''070196)\cos(5n^{\mathrm{V}}t - 2n^{\mathrm{IV}}t + 5\varepsilon^{\mathrm{V}} - 2\varepsilon^{\mathrm{IV}})
\end{aligned}
\]
Hiedurch wird der obige Fehler in der Länge auf $5'\,22''$ reducirt, also so klein, dass er ohne Bedenken der Beobachtung zugeschrieben werden darf.

\clearpage
\fancyhead[L]{\small HUMBOLDT. CONSPECTUS LONGITUDINUM ETC.}
\hypertarget{native-514-lower}{}\label{native-514-lower}
\begin{center}\small Göttingische gelehrte Anzeigen. Stück 193. Seite 1928. 1808 December 3.\\[3pt]\rule{.72\linewidth}{.3pt}\end{center}

\emph{Conspectus longitudinum et latitudinum geographicarum, per decursum annorum} 1799 \emph{ad} 1804 \emph{in plaga aequinoctiali ab} \textsc{Alex. de Humboldt} \emph{astronomice observatarum, calculo subjecit} \textsc{Jabbo Oltmans}. 16 Seiten in Folio. Bei \textsc{F. Schöll} und \textsc{J. G. Cotta} in Paris und Tübingen.

Der kostbaren Ausbeute, welche die Geographie der Amerikanischen Tropenländer durch die Reise des Hrn. \textsc{von Humboldt} erhalten hat, wird bekanntlich eine eigne Abtheilung seines grossen Werks gewidmet, worin wir nicht blos alle seine zahlreichen astronomischen Beobachtungen in extenso, sondern zugleich eine sehr sorgfältige Discussion derselben von dem bereits vortheilhaft bekannten Hrn. \textsc{Jabbo Oltmans} zu erwarten haben. Da indess der Druck dieses Werks eine beträchtliche Zeit erfordern wird, so ist einstweilen gegenwärtige summarische Zusammenstellung der vornehmsten Ortsbestimmungen veranstaltet worden, welche blos die letzten Resultate enthält. Nahe an dreihundert Bestimmungen werden hier geliefert, die grösstentheils auf Hrn. \textsc{von Humboldts} eigne Beobachtungen gegründet sind. Man kann nicht umhin, die Thätigkeit und Vielseitigkeit des Mannes zu bewundern, der so viele verschiedenartige Zwecke zugleich umfassend, jeden so zu erreichen wusste, als wenn er ihm seine Zeit und Kräfte ungetheilt hätte widmen können.
